\section{Planteamiento del problema}

En Colombia, la medición de la pobreza es un insumo crítico para el diseño, seguimiento y evaluación de las políticas públicas de bienestar social y equidad territorial. Este análisis se aborda mediante dos enfoques complementarios: la pobreza monetaria, que identifica a los hogares con ingresos insuficientes para adquirir una canasta básica de bienes y servicios, y la pobreza multidimensional, instrumentada a través del Índice de Pobreza Multidimensional (IPM), que evalúa privaciones estructurales en educación, salud, trabajo, niñez y vivienda. Mientras el primer enfoque examina la suficiencia de ingresos, el IPM captura las carencias materiales y de capacidades que afectan el bienestar de los hogares \cite{DANE_PobrezaMultidimensional_2023}.

No obstante, las estimaciones oficiales del Departamento Administrativo Nacional de Estadística (DANE) se desagregan, por lo general, a nivel departamental o para las principales áreas metropolitanas. Esta escala analítica diluye las heterogeneidades locales y oculta los contrastes socioeconómicos existentes entre municipios \cite{DANE_PobrezaMultidimensional_2023}. La ausencia de datos confiables a escala municipal debilita la focalización de programas sociales, la distribución eficiente de recursos y la formulación de políticas adaptadas a las realidades de las regiones históricamente rezagadas. La raíz del problema es metodológica: las encuestas de hogares empleadas para estas mediciones —como la Gran Encuesta Integrada de Hogares (GEIH) o la Encuesta Nacional de Calidad de Vida (ECV)— poseen diseños muestrales representativos a nivel departamental, pero carecen de la potencia estadística necesaria para realizar inferencias válidas en el ámbito municipal \cite{DANE_ENCV_2022}.

A esta restricción espacial se suma una limitación temporal. Colombia carece de mediciones oficiales directas y frecuentes del IPM municipal, lo que impide monitorear oportunamente la evolución de las privaciones en unidades territoriales pequeñas. Las estadísticas oficiales con mayor desagregación geográfica provienen de los censos de población y vivienda debido a su cobertura exhaustiva. Sin embargo, al ser operaciones periódicas y no anuales, los censos no permiten un seguimiento continuo. Este desfase intercensal genera un vacío de información que restringe la actualización del IPM municipal y reduce la capacidad institucional para detectar cambios recientes en las condiciones de vida de la población \cite{DANE_PobrezaMultidimensional_2023} \cite{DANE_ENCV_2022}.

Ante la inviabilidad de obtener mediciones directas actualizadas a nivel municipal, este proyecto implementa la estimación indirecta del IPM. En un entorno de alta desigualdad y heterogeneidad territorial como el colombiano, los promedios departamentales suelen enmascarar brechas críticas; dentro de un mismo departamento coexisten municipios con estándares de bienestar urbanos junto a otros con niveles de pobreza propios de las zonas rurales más vulnerables \cite{CEPAL_PanoramaSocial_2023}. Para resolver esta restricción, la estadística aplicada y la ciencia de datos ofrecen metodologías de Estimación para Áreas Pequeñas (Small Area Estimation, SAE), diseñadas para producir inferencias precisas en dominios con muestras insuficientes mediante la integración de múltiples fuentes de información \cite{RaoMolina_2015}.

Estos enfoques combinan los datos de encuestas con variables auxiliares procedentes de censos, registros administrativos o imágenes satelitales, empleando modelos jerárquicos o mixtos que incorporan estructuras espaciales \cite{MolinaRao_2010}. Al aprovechar la correlación geográfica entre unidades territoriales contiguas, estos modelos mejoran la precisión y la coherencia espacial de las estimaciones \cite{PratesiSalvati_2019}. En el contexto nacional, el uso de metodologías SAE se encuentra en desarrollo. Aunque estudios recientes demuestran su potencial para optimizar la medición de indicadores sociales a escala subnacional \cite{DANE_PobrezaMultidimensional_2023} \cite{LopezSanchezPerez_2021}, su aplicación al IPM sigue siendo incipiente. La adopción de modelos mixtos con componentes espaciales representa un avance metodológico clave para generar estimaciones consistentes, especialmente en municipios con baja representatividad muestral o registros administrativos deficientes.

La resolución de este problema técnico es indispensable para consolidar una gestión pública basada en evidencia, donde la disponibilidad de datos robustos a nivel municipal es el eje para optimizar la asignación presupuestal y formular políticas con precisión geoespacial \cite{PNUD_HDR_2022}. Desde la perspectiva de la ciencia de datos, el desafío radica en el desarrollo de arquitecturas analíticas capaces de integrar fuentes heterogéneas para mitigar la invisibilidad estadística de las comunidades más vulnerables, cuya situación socioeconómica suele quedar subrepresentada en los agregados oficiales.